% !TeX root = ../../../../../main.tex
% chapters/sections/chapter5/subsections/a_shock/subsubsections/interpretations/overview.tex

\subsubsection{2平面分析によるグラフの解釈の概要}
\label{sec:chapter5_a_shock_interpretations_overview}
ここでは\(a\)ショックが発生してからの主要変数の動きについて、
名目利子率\(i_t^H\)の反応が異なる2つの政策群に分けて、2平面分析による解釈をおこなう。
グラフの解釈をおこなうにあたり\(a\)ショックの際に各曲線がもつ性質をまとめておく。
\begin{itemize}
    \item YS-\(\pi\)曲線（差分形式）の右辺の（定数倍を除く）変動要因：
           \(\hat{y}_t^H-\hat{a}_t^H\)
    \item YS-\(\hat{p}\),CS-\(\hat{p}\)曲線の切片：
          \(\hat{p}_{t-1}^H+\beta_{ss}^H\operatorname{E}_t[\hat{\pi}_{t+1}^H]-2\kappa^H\hat{a}_t^H\)
    \item YS-\(\hat{p}\)の傾き\(<\)CS-\(\hat{p}\)曲線の傾き（\(0<\alpha^H<1\)より）：
          \(2\kappa^H<\frac{2\kappa^H}{\alpha^H}+\frac{1-\alpha^H}{\alpha^H}\) 
    \item YS-\(\hat{p}\)曲線と完全伸縮価格\(\xi^H\to0\)による垂直なYS-\(\hat{p}\)曲線の\(\hat{y}_t^H\)切片：
          \(\hat{a}_t^H\)
    \item YD-\(\hat{p}\),CD-\(\hat{p}\)曲線の切片：
          \(-\hat{\lambda}_t^H\)
\end{itemize}